% CERP Exam Question Bank - Set 4 (Questions 601-800)
% Additional 200 Unique Questions - Not Repeated from Sets 1-3
% Compiled from ABA CERP Exam Content Outline
% Compiled on: 2026-06-13

\documentclass[12pt, letterpaper]{article}
\usepackage[utf8]{inputenc}
\usepackage{geometry}
\usepackage{amsmath}
\usepackage{amsfonts}
\usepackage{amssymb}
\usepackage{titlesec}
\usepackage{enumitem}
\usepackage{xcolor}
\usepackage{fancyhdr}
\usepackage{multicol}

\geometry{margin=1in}
\setlength{\parindent}{0pt}
\setlength{\parskip}{0.8ex}

\pagestyle{fancy}
\fancyhf{}
\fancyhead[L]{\textbf{CERP Exam Prep - Set 4}}
\fancyhead[R]{\thepage}
\fancyfoot[C]{© 2026 Enterprise Risk Training}

\newcounter{questioncounter}
\setcounter{questioncounter}{601}

\newcommand{\question}{
    \par\noindent
    \textbf{\arabic{questioncounter}.}\quad
    \stepcounter{questioncounter}
}

\newcommand{\answer}[1]{
    \par\noindent
    \textbf{Answer:} #1
    \vspace{0.3cm}
    \par
}

\titleformat{\section}{\large\bfseries}{\thesection}{1em}{}
\titleformat{\subsection}{\normalsize\bfseries}{\thesubsection}{1em}{}

\begin{document}

\begin{center}
    \vspace*{1cm}
    {\huge \textbf{Certified Enterprise Risk Professional (CERP)}} \\[0.5cm]
    {\huge \textbf{Exam Question Bank - Set 4}} \\[1cm]
    {\Large 200 Additional Multiple-Choice Questions (601-800)} \\[1.5cm]
    {\large Based on ABA CERP Exam Content Outline} \\[0.3cm]
    {\large Domains 1-8: Risk Governance \& Risk Management} \\[2cm]
    {\large \textbf{Questions 601 - 800}} \\[0.3cm]
    \vfill
    {\large \today}
\end{center}

\newpage

% ============================================================
% DOMAIN 1: Risk Governance (Questions 601-616)
% ============================================================
\section{RISK GOVERNANCE}

\subsection{Domain 1: Board and Senior Management Oversight (Questions 601-616)}

\question What is the primary mechanism for board oversight of risk?
A) Daily transaction approval
B) Regular receipt and review of risk reports and committee minutes
C) Trading on the floor
D) Approving every loan
\answer{B}

\question Which is a best practice for board risk committee composition?
A) All executives
B) Majority independent directors with relevant expertise
C) All operations staff
D) External consultants only
\answer{B}

\question What is the board's role in risk culture assessment?
A) No role
B) Periodically assess whether risk culture aligns with stated values
C) Delegate to junior staff
D) Ignore culture
\answer{B}

\question How should management report emerging risks to the board?
A) Never report
B) Proactively identify and report emerging trends with potential impact
C) Only report after loss
D) Only report positive news
\answer{B}

\question What is the consequence of a board that does not understand the bank's risk profile?
A) Strong oversight
B) Inability to provide effective challenge and oversight
C) Higher profitability
D) Lower regulatory scrutiny
\answer{B}

\question Which is a key responsibility of the board's audit committee?
A) Set risk appetite
B) Oversee financial reporting, internal controls, and audit functions
C) Approve individual loans
D) Manage daily operations
\answer{B}

\question What is the role of "board minutes" in risk governance?
A) No purpose
B) Document decisions, challenges, and oversight activities
C) Marketing material
D) Employee evaluations
\answer{B}

\question How does the board ensure risk appetite is operationalized?
A) By ignoring implementation
B) By requiring management to embed appetite into limits and reporting
C) By setting appetite once and forgetting it
D) By delegating entirely
\answer{B}

\question Which is a sign of ineffective board risk oversight?
A) Frequent risk discussions
B) Management always agrees with board; no credible challenge
C) Independent risk reports
D) Diverse board expertise
\answer{B}

\question What is the board's role in crisis management?
A) No role
B) Oversee crisis response plans and ensure recovery
C) Execute the crisis plan
D) Ignore the crisis
\answer{B}

\question Which should be included in the board's risk information package?
A) Only financial results
B) Key risk indicators, limit breaches, emerging risks, and issue status
C) Employee birthdays
D) Office supply orders
\answer{B}

\question What is the board's responsibility regarding whistleblower reports?
A) Ignore them
B) Ensure proper investigation and non-retaliation
C) Punish whistleblowers
D) Delegate without oversight
\answer{B}

\question How often should the board review the enterprise risk management framework?
A) Never
B) At least annually
C) Once every 5 years
D) Only after a loss
\answer{B}

\question Which is a key question the board should ask management about risk?
A) "What is the weather forecast?"
B) "What are our top three risks and what are we doing about them?"
C) "What is for lunch?"
D) "Who is the newest employee?"
\answer{B}

\question What is the board's role in risk appetite breaches?
A) Ignore breaches
B) Review material breaches and management's remediation plan
C) Approve all breaches automatically
D) Delegate to lowest level
\answer{B}

\question Which is a key attribute of an effective board risk committee?
A) No independent members
B) Access to independent expertise and direct communication with risk management
C) No meeting minutes
D) Infrequent meetings
\answer{B}

% ============================================================
% DOMAIN 2: Policies, Procedures and Limits (Questions 617-640)
% ============================================================

\subsection{Domain 2: Policies, Procedures and Limits (Questions 617-640)}

\question What is a "limit monitoring system"?
A) A suggestion
B) A system that tracks exposures against limits in near real-time
C) A marketing system
D) An HR system
\answer{B}

\question Who is responsible for establishing the risk limit framework?
A) Front-line staff
B) Board of Directors or delegated committee
C) External auditors
D) Customers
\answer{B}

\question What is a "limit breach investigation"?
A) No action
B) A review to determine cause, impact, and remediation
C) Automatic approval
D) Ignoring the breach
\answer{B}

\question Which is a typical limit for market risk in a trading book?
A) Loan-to-value ratio
B) Maximum daily Value at Risk (VaR) limit
C) Employee turnover
D) Customer satisfaction
\answer{B}

\question What is a "policy attestation"?
A) No confirmation
B) Employee confirmation that they have read and understand policies
C) A marketing campaign
D) An HR form
\answer{B}

\question Which is a regulatory expectation for model risk limits?
A) No limits
B) Limits on model use and thresholds for model performance metrics
C) Unlimited model use
D) No performance monitoring
\answer{B}

\question What is a "limit breach escalation matrix"?
A) No matrix
B) A document defining who to notify based on breach size and severity
C) A marketing tool
D) An employee directory
\answer{B}

\question Which is a typical policy for gifts and entertainment?
A) Unlimited gifts
B) Prohibition or limits on gifts to/from vendors and disclosure requirements
C) No policy
D) Encouraging gifts
\answer{B}

\question What is a "limit waiver"?
A) Permanent elimination of limit
B) Temporary, documented permission to exceed a limit
C) An automatic approval
D) A marketing approval
\answer{B}

\question Which is a regulatory expectation for anti-money laundering (AML) policies?
A) No policy
B) Customer due diligence, transaction monitoring, and SAR filing
C) Ignore suspicious activity
D) No training
\answer{B}

\question What is a "policy breach root cause analysis"?
A) Ignoring the cause
B) Investigation to identify why the breach occurred
C) Blaming someone
D) No analysis
\answer{B}

\question Which is a typical limit for foreign exchange (FX) risk?
A) Loan-to-value
B) Maximum open position limit
C) Employee count
D) Number of branches
\answer{B}

\question What is a "policy review cycle"?
A) No review
B) Scheduled periodic review of policies for relevance and effectiveness
C) One-time review
D) Review only after loss
\answer{B}

\question Which is a regulatory expectation for data privacy policies (e.g., GDPR)?
A) No policy
B) Policies for data collection, consent, access, and breach notification
C) Ignore customer data
D) Unlimited data sharing
\answer{B}

\question What is a "limit breach remediation"?
A) No action
B) Actions to address the breach and prevent recurrence
C) Approving the breach
D) Ignoring the breach
\answer{B}

\question Which is a typical policy for political contributions?
A) Unlimited contributions
B) Approval requirements and disclosure
C) No policy
D) Encouraging contributions
\answer{B}

\question What is a "risk limit exception report"?
A) No report
B) A report listing all limit exceptions with approval and rationale
C) A marketing report
D) An HR report
\answer{B}

\question Which is a regulatory expectation for consumer protection policies?
A) No policy
B) Fair lending, UDAAP prohibitions, and complaint handling
C) Ignore consumer complaints
D) No disclosure requirements
\answer{B}

\question What is a "policy training program"?
A) No training
B) Training to ensure employees understand and can comply with policies
C) Optional training
D) Training only for executives
\answer{B}

\question Which is a typical limit for commodities risk?
A) Loan-to-value
B) Maximum position limit by commodity type
C) Employee count
D) Branch count
\answer{B}

\question What is a "policy exception trend analysis"?
A) No analysis
B) Analyzing exception patterns to identify control gaps
C) Ignoring trends
D) Only looking at one exception
\answer{B}

\question Which is a regulatory expectation for fair lending policies?
A) No policy
B) Policies prohibiting discrimination and requiring monitoring
C) Encouraging discrimination
D) No data collection
\answer{B}

\question What is a "limit breach dashboard"?
A) No dashboard
B) Visual display of breaches, severity, and remediation status
C) Marketing dashboard
D) HR dashboard
\answer{B}

\question Which is a typical policy for records retention?
A) Keep everything forever
B) Retention schedules based on legal and regulatory requirements
C) Delete everything immediately
D) No retention policy
\answer{B}

% ============================================================
% DOMAIN 3: MIS and Data Governance (Questions 641-655)
% ============================================================

\subsection{Domain 3: Management Information Systems (Questions 641-655)}

\question What is "data stewardship"?
A) No role
B) Day-to-day management of data quality and metadata
C) Deleting data
D) Marketing data
\answer{B}

\question Which is a risk of multiple data sources for the same metric?
A) Consistency
B) Data reconciliation challenges and potential inconsistencies
C) Perfect accuracy
D) Lower cost
\answer{B}

\question What is a "data quality issue log"?
A) No log
B) A log to track data quality problems and resolution
C) A marketing log
D) An HR log
\answer{B}

\question Which is a key element of a data governance operating model?
A) No roles
B) Defined roles (data owner, steward, custodian) and responsibilities
C) Everyone does everything
D) No accountability
\answer{B}

\question What is a "data quality scorecard"?
A) No scorecard
B) Metrics and KPIs measuring data quality dimensions
C) Marketing scorecard
D) Employee scorecard
\answer{B}

\question Which is a control for data lineage documentation?
A) No documentation
B) Documenting data flows from source to report
C) Ignoring data sources
D) No traceability
\answer{B}

\question What is a "data retention schedule"?
A) Keep all data forever
B) Defined time periods for retaining different data types
C) Delete all data immediately
D) No schedule
\answer{B}

\question Which is a risk of manual data reconciliation?
A) Perfect accuracy
B) Human error and time delays
C) Lower cost
D) Faster processing
\answer{B}

\question What is a "data quality metric"?
A) No metric
B) Quantitative measure of data accuracy, completeness, or timeliness
C) Marketing metric
D) HR metric
\answer{B}

\question Which is a best practice for risk data architecture?
A) Siloed spreadsheets
B) Enterprise data warehouse with defined data models
C) No integration
D) Manual entry only
\answer{B}

\question What is a "data validation rule set"?
A) No rules
B) Collection of automated checks for data quality
C) Marketing rules
D) Employee rules
\answer{B}

\question Which is a control for data privacy?
A) No controls
B) Data masking, access controls, and retention limits
C) Unlimited access to all data
D) No encryption
\answer{B}

\question What is a "data quality dashboard" used for?
A) No purpose
B) To monitor and report data quality metrics to stakeholders
C) Marketing only
D) HR only
\answer{B}

\question Which is a key principle of data governance?
A) No accountability
B) Data is an asset with defined ownership
C) Anyone can change data
D) No documentation
\answer{B}

\question What is a "data issue escalation process"?
A) No process
B) Defined steps for escalating data quality problems
C) Ignoring issues
D) Only manual notification
\answer{B}

% ============================================================
% DOMAIN 4: Control Framework (Questions 656-670)
% ============================================================

\subsection{Domain 4: Control Framework (Questions 656-670)}

\question What is a "control deficiency remediation plan"?
A) No plan
B) Documented actions to fix a control weakness
C) Ignoring the deficiency
D) Marketing plan
\answer{B}

\question Which is a control for physical security?
A) No controls
B) Access badges, cameras, and security guards
C) Unlimited access
D) No monitoring
\answer{B}

\question What is a "control self-assessment" (CSA) used for?
A) Eliminating controls
B) Empowering process owners to evaluate control effectiveness
C) Outsourcing control testing
D) Ignoring controls
\answer{B}

\question Which is a control for IT change management?
A) No controls
B) Segregation of duties between development, testing, and production access
C) Anyone can change production
D) No documentation
\answer{B}

\question What is a "control testing plan"?
A) No plan
B) Documented approach for testing control design and operation
C) A marketing plan
D) An HR plan
\answer{B}

\question Which is a control for vendor access to bank systems?
A) No controls
B) Access reviews, contracts, and monitoring
C) Unlimited vendor access
D) No contracts
\answer{B}

\question What is a "control effectiveness rating"?
A) No rating
B) Assessment of whether control is designed and operating effectively
C) Marketing rating
D) Employee rating
\answer{B}

\question Which is a control for data backup?
A) No backup
B) Regular backups stored offsite with testing
C) No recovery testing
D) Single copy only
\answer{B}

\question What is a "control gap analysis"?
A) No analysis
B) Identifying risks without adequate controls
C) Ignoring gaps
D) Only reviewing existing controls
\answer{B}

\question Which is a control for insider threat?
A) No controls
B) User activity monitoring and segregation of duties
C) Unlimited access for all
D) No logging
\answer{B}

\question What is a "control rationalization"?
A) Adding unnecessary controls
B) Reviewing and removing redundant or ineffective controls
C) Ignoring control efficiency
D) No review
\answer{B}

\question Which is a control for financial close process?
A) No controls
B) Reconciliation, review, and approval of journal entries
C) Unlimited adjustments
D) No approval
\answer{B}

\question What is a "control documentation standard"?
A) No standard
B) Consistent template for documenting control attributes
C) Any format allowed
D) Verbal documentation only
\answer{B}

\question Which is a control for privileged access?
A) No controls
B) Approval, logging, and periodic review of privileged accounts
C) Unlimited privileged access
D) No approval
\answer{B}

\question What is a "control monitoring alert"?
A) No alert
B) Automated notification when a control fails or detects an anomaly
C) Marketing alert
D) HR alert
\answer{B}

% ============================================================
% DOMAIN 5: Risk Identification (Questions 671-695)
% ============================================================

\subsection{Domain 5: Risk Identification (Questions 671-695)}

\question Which is a risk from intellectual property (IP) infringement?
A) No risk
B) Legal liability and reputational damage
C) Guaranteed profit
D) Lower costs
\answer{B}

\question What is "scenario-based risk identification"?
A) No scenarios
B) Using hypothetical scenarios to identify potential risks
C) Only historical data
D) Ignoring future risks
\answer{B}

\question Which is a risk from employee misconduct?
A) No risk
B) Legal, regulatory, and reputational consequences
C) Guaranteed compliance
D) Higher morale
\answer{B}

\question What is "risk identification frequency" based on?
A) Random
B) Risk velocity and business change (e.g., annual, quarterly, continuous)
C) Once only
D) Never
\answer{B}

\question Which is a risk from inadequate documentation?
A) No risk
B) Operational inefficiency and audit findings
C) Guaranteed accuracy
D) Lower costs
\answer{B}

\question What is a "risk identification workshop"?
A) A party
B) Facilitated session with stakeholders to brainstorm risks
C) A solo activity
D) A marketing event
\answer{B}

\question Which is a risk from reliance on manual processes?
A) No risk
B) Higher error rates and slower response times
C) Perfect accuracy
D) Lower cost
\answer{B}

\question What is "industry benchmarking" for risk identification?
A) Ignoring peers
B) Comparing risk profiles and incidents with peer institutions
C) Copying competitors blindly
D) No comparison
\answer{B}

\question Which is a risk from supply chain disruption?
A) No risk
B) Inability to deliver services or products to customers
C) Guaranteed resilience
D) Lower costs
\answer{B}

\question What is a "risk identification interview"?
A) A job interview
B) One-on-one discussion with process owners to surface risks
C) A marketing call
D) An HR meeting
\answer{B}

\question Which is a risk from product complexity?
A) No risk
B) Higher likelihood of errors and poor customer understanding
C) Guaranteed sales
D) Lower cost
\answer{B}

\question What is "external loss data" used for?
A) No use
B) To identify risks that have affected other organizations
C) To ignore industry events
D) Marketing only
\answer{B}

\question Which is a risk from key person dependency?
A) No risk
B) Operational disruption if a critical employee leaves
C) Guaranteed continuity
D) Lower cost
\answer{B}

\question What is a "risk identification questionnaire"?
A) No questionnaire
B) Structured survey to gather risk information from staff
C) Marketing survey
D) Employee satisfaction survey
\answer{B}

\question Which is a risk from rapid scaling of operations?
A) No risk
B) Control breakdowns and increased error rates
C) Guaranteed success
D) Lower costs
\answer{B}

\question What is "regulatory change monitoring" for risk identification?
A) Ignore regulations
B) Tracking proposed and final regulations to identify new compliance risks
C) Only after enforcement action
D) No monitoring
\answer{B}

\question Which is a risk from business model disruption?
A) No risk
B) Loss of revenue from new competitors or technology
C) Guaranteed market share
D) Lower costs
\answer{B}

\question What is a "risk identification trigger event"?
A) Routine activity
B) A specific occurrence that prompts a risk reassessment
C) A marketing campaign
D) An employee meeting
\answer{B}

\question Which is a risk from poor data quality?
A) No risk
B) Misinformed decisions and regulatory fines
C) Perfect decisions
D) Lower costs
\answer{B}

\question What is "lessons learned" for risk identification?
A) No value
B) Using past incidents and near-misses to identify future risks
C) Ignoring history
D) Only focusing on successes
\answer{B}

\question Which is a risk from geographic concentration?
A) No risk
B) Exposure to localized economic or natural disaster events
C) Diversification benefit
D) Lower risk
\answer{B}

\question What is a "risk identification taxonomy"?
A) No classification
B) Hierarchical structure of risk categories for consistent identification
C) Marketing classification
D) Employee classification
\answer{B}

\question Which is a risk from conflict of interest?
A) No risk
B) Unfair decisions and regulatory violations
C) Guaranteed fairness
D) Lower costs
\answer{B}

\question What is "red teaming" for risk identification?
A) No technique
B) Adversarial challenge to identify blind spots
C) Groupthink
D) Ignoring challenges
\answer{B}

\question Which is a risk from inadequate incident response planning?
A) No risk
B) Delayed recovery and amplified losses
C) Guaranteed recovery
D) Lower costs
\answer{B}

% ============================================================
% DOMAIN 6: Risk Measurement (Questions 696-720)
% ============================================================

\subsection{Domain 6: Risk Measurement and Evaluation (Questions 696-720)}

\question What is "risk contribution"?
A) Total risk
B) The incremental risk added by a specific business unit or position
C) Average risk
D) Minimum risk
\answer{B}

\question Which is a measure of spread risk in credit?
A) Duration
B) Credit spread widening impact on portfolio value
C) Loan-to-value
D) Debt-to-income
\answer{B}

\question What is "prepayment risk"?
A) No risk
B) Risk that borrowers repay loans earlier than expected
C) Default risk
D) Interest rate risk
\answer{B}

\question Which is a measure of option risk?
A) Delta, Gamma, Vega, Theta
B) Loan-to-value
C) Debt-to-income
D) Employee count
\answer{A}

\question What is "liquidity at risk" (LaR)?
A) No measure
B) Potential liquidity shortfall under stress scenarios
C) Profit measure
D) Capital measure
\answer{B}

\question Which is a measure of basis risk in commodities?
A) Price correlation
B) Difference between local and benchmark prices
C) Volatility
D) Default probability
\answer{B}

\question What is "recovery rate" (RR) in credit risk?
A) 1 - LGD
B) Default probability
C) Exposure amount
D) Loss amount
\answer{A}

\question Which is a measure of counterparty risk exposure?
A) Loan-to-value
B) Current Exposure (CE) plus Potential Future Exposure (PFE)
C) Debt-to-income
D) Employee count
\answer{B}

\question What is "wrong-way risk" measurement?
A) No measurement
B) Correlation between exposure and counterparty credit quality
C) Independent risk
D) Operational risk
\answer{B}

\question Which is a measure of model risk?
A) Model validation status and performance metrics
B) Employee count
C) Branch count
D) Loan balance
\answer{A}

\question What is "risk appetite statement" quantification?
A) No quantification
B) Specific metrics, thresholds, and limits tied to risk appetite
C) Only qualitative
D) Marketing metrics
\answer{B}

\question Which is a measure of earnings volatility?
A) Standard deviation of net income
B) Loan balance
C) Deposit balance
D) Employee count
\answer{A}

\question What is "credit conversion factor" (CCF)?
A) No factor
B) Percentage of off-balance sheet exposure converted to credit equivalent
C) Default probability
D) Loss given default
\answer{B}

\question Which is a measure of operational risk capital?
A) Basic Indicator Approach or Standardized Approach
B) Loan-to-value
C) Debt-to-income
D) Duration
\answer{A}

\question What is "risk appetite utilization report"?
A) No report
B) Current risk levels compared to appetite limits
C) Marketing report
D) HR report
\answer{B}

\question Which is a measure of foreign exchange risk?
A) FX VaR and sensitivity (delta)
B) Loan-to-value
C) Debt-to-income
D) Employee count
\answer{A}

\question What is "downgrade risk" in credit?
A) No risk
B) Impact of credit rating downgrade on portfolio value
C) Default risk only
D) Prepayment risk
\answer{B}

\question Which is a measure of concentration risk in loan portfolio?
A) Herfindahl-Hirschman Index (HHI)
B) Total loans
C) Average loan size
D) Number of loans
\answer{A}

\question What is "liquidity gap analysis"?
A) No analysis
B) Cumulative net cash flow mismatch over time buckets
C) Profit analysis
D) Capital analysis
\answer{B}

\question Which is a measure of interest rate risk sensitivity?
A) DV01 (dollar value of 1 basis point)
B) Loan balance
C) Deposit balance
D) Employee count
\answer{A}

\question What is "stress test scenario severity"?
A) No severity
B) Magnitude of hypothetical adverse conditions
C) Profit level
D) Employee count
\answer{B}

\question Which is a measure of basis risk in interest rates?
A) Correlation between different rate indices
B) Loan-to-value
C) Debt-to-income
D) Employee count
\answer{A}

\question What is "risk measurement validation"?
A) No validation
B) Independent review of risk models and calculations
C) Trusting without review
D) Marketing validation
\answer{B}

\question Which is a measure of rollover risk?
A) Percentage of liabilities maturing within 30 days
B) Loan balance
C) Deposit balance
D) Employee count
\answer{A}

\question What is "risk measurement frequency"?
A) Once only
B) Based on risk velocity (daily, weekly, monthly, quarterly)
C) Never
D) Only at year-end
\answer{B}

% ============================================================
% DOMAIN 7: Risk Responses (Questions 721-745)
% ============================================================

\subsection{Domain 7: Risk Responses (Questions 721-745)}

\question When is "risk avoidance" the best response?
A) Always
B) When risk cannot be mitigated to within appetite cost-effectively
C) Never
D) For all operational risks
\answer{B}

\question What is "risk mitigation prioritization"?
A) No prioritization
B) Focusing resources on highest residual risks first
C) Random selection
D) Mitigating all risks equally
\answer{B}

\question Which is an example of risk transfer via securitization?
A) No transfer
B) Selling loan portfolios to remove credit risk
C) Accepting all credit risk
D) Avoiding lending
\answer{B}

\question What is a "risk response cost-benefit analysis"?
A) No analysis
B) Comparing cost of response to expected risk reduction benefit
C) Always choose most expensive
D) Never analyze
\answer{B}

\question Which is an example of risk mitigation for business continuity?
A) No plan
B) Redundant systems and alternate sites
C) Accepting all disruption
D) Avoiding all operations
\answer{B}

\question What is "risk response tracking"?
A) No tracking
B) Monitoring implementation and effectiveness of risk responses
C) Ignoring responses
D) Only planning
\answer{B}

\question Which is an example of risk acceptance for strategic risk?
A) No strategy
B) Pursuing a high-reward strategy knowingly within appetite
C) Avoiding all strategy
D) Transferring strategy risk
\answer{B}

\question What is a "risk response action plan"?
A) No plan
B) Documented steps, owners, and timelines for response implementation
C) Marketing plan
D) HR plan
\answer{B}

\question Which is an example of risk mitigation for legal risk?
A) No legal review
B) Contract reviews and legal advice before commitments
C) Accepting all legal liability
D) Avoiding contracts
\answer{B}

\question What is "risk response effectiveness testing"?
A) No testing
B) Verifying that risk responses actually reduce risk as intended
C) Assuming effectiveness
D) No verification
\answer{B}

\question Which is an example of risk transfer via outsourcing?
A) No transfer
B) Contractually shifting operational risk to vendor
C) Keeping all operational risk
D) Avoiding outsourcing
\answer{B}

\question What is a "risk response register"?
A) No register
B) Log of all risk responses with status and owners
C) Marketing log
D) Employee log
\answer{B}

\question Which is an example of risk mitigation for compliance risk?
A) No compliance program
B) Automated compliance monitoring systems
C) Accepting all violations
D) Avoiding all business
\answer{B}

\question What is "residual risk reporting"?
A) No reporting
B) Reporting risk remaining after responses to management
C) Only inherent risk
D) Only control risk
\answer{B}

\question Which is an example of risk avoidance for country risk?
A) Lending in high-risk countries
B) Ceasing operations in a high-risk country
C) Accepting expropriation risk
D) No action
\answer{B}

\question What is "risk response escalation"?
A) No escalation
B) Raising risk response decisions to appropriate authority levels
C) Making all decisions at low level
D) Ignoring authority
\answer{B}

\question Which is an example of risk mitigation for data privacy risk?
A) No controls
B) Data encryption and access controls
C) Accepting all breaches
D) Avoiding data collection
\answer{B}

\question What is "risk response validation"?
A) No validation
B) Independent confirmation that response is effective
C) Self-validation only
D) No confirmation
\answer{B}

\question Which is an example of risk transfer via derivatives?
A) No transfer
B) Using interest rate swaps to hedge fixed-rate loan exposure
C) Accepting all rate risk
D) Avoiding derivatives
\answer{B}

\question What is "risk response ownership"?
A) No owner
B) Assigned person responsible for implementing and maintaining response
C) Everyone owns
D) No one owns
\answer{B}

\question Which is an example of risk mitigation for reputational risk?
A) No controls
B) Customer complaint handling and social media monitoring
C) Accepting all reputational damage
D) Avoiding public interaction
\answer{B}

\question What is "risk response documentation"?
A) No documentation
B) Written evidence of response decisions and implementation
C) Verbal only
D) No evidence
\answer{B}

\question Which is an example of risk avoidance for concentration risk?
A) No diversification
B) Capping exposure to any single counterparty
C) Accepting concentration
D) No limit
\answer{B}

\question What is "risk response monitoring"?
A) No monitoring
B) Ongoing review of response effectiveness and changes in risk
C) One-time review
D) No review
\answer{B}

\question Which is an example of risk mitigation for third-party risk?
A) No vendor management
B) Vendor due diligence and performance scorecards
C) Accepting all vendor failures
D) Avoiding all vendors
\answer{B}

% ============================================================
% DOMAIN 8: Risk Monitoring (Questions 746-770)
% ============================================================

\subsection{Domain 8: Risk Monitoring (Questions 746-770)}

\question What is a "control monitoring frequency"?
A) One-time
B) Risk-based schedule for testing control effectiveness
C) Never
D) Only after loss
\answer{B}

\question Which is a KRI for project risk?
A) Loan growth
B) Number of missed milestones or budget overruns
C) Deposit balance
D) Employee count
\answer{B}

\question What is "risk reporting distribution list"?
A) No list
B) Defined recipients for each risk report
C) Send to everyone
D) Send to no one
\answer{B}

\question Which is a KRI for system availability?
A) Loan balance
B) System uptime percentage and number of outages
C) Deposit balance
D) Employee count
\answer{B}

\question What is "risk monitoring threshold"?
A) A suggestion
B) A pre-defined level that triggers review or action
C) A final loss
D) A marketing target
\answer{B}

\question Which is a KRI for customer satisfaction risk?
A) Loan growth
B) Net Promoter Score (NPS) or complaint volume
C) Deposit balance
D) Capital ratio
\answer{B}

\question What is "risk escalation protocol"?
A) No protocol
B) Defined process for raising risk issues to higher authority
C) Ignoring escalation
D) Only verbal escalation
\answer{B}

\question Which is a KRI for employee turnover risk?
A) Loan balance
B) Voluntary turnover rate in key risk roles
C) Deposit balance
D) Branch count
\answer{B}

\question What is "risk reporting automation"?
A) Manual reports only
B) Automated generation and distribution of risk reports
C) No reporting
D) Spreadsheet only
\answer{B}

\question Which is a KRI for training risk?
A) Loan growth
B) Percentage of employees with required training complete
C) Deposit balance
D) Capital ratio
\answer{B}

\question What is "risk monitoring committee"?
A) No committee
B) Regular meeting to review risk metrics and breaches
C) Marketing committee
D) HR committee
\answer{B}

\question Which is a KRI for data quality risk?
A) Loan balance
B) Error rates in critical risk reports
C) Deposit balance
D) Employee count
\answer{B}

\question What is "risk reporting retention"?
A) No retention
B) Keeping risk reports per regulatory and audit requirements
C) Delete all reports immediately
D) Keep forever
\answer{B}

\question Which is a KRI for vendor performance?
A) Loan growth
B) Vendor service level agreement (SLA) compliance percentage
C) Deposit balance
D) Employee count
\answer{B}

\question What is "risk monitoring system"?
A) No system
B) Technology platform for tracking risk metrics and limits
C) Manual process only
D) Spreadsheet only
\answer{B}

\question Which is a KRI for change management?
A) Loan balance
B) Number of emergency changes without proper approval
C) Deposit balance
D) Employee count
\answer{B}

\question What is "risk reporting remediation tracking"?
A) No tracking
B) Monitoring of actions to address report findings
C) Ignoring findings
D) Only reporting
\answer{B}

\question Which is a KRI for access control risk?
A) Loan growth
B) Number of users with excessive or stale privileges
C) Deposit balance
D) Employee count
\answer{B}

\question What is "risk monitoring continuous"?
A) Periodic only
B) Real-time or near-real-time risk surveillance
C) Manual only
D) No monitoring
\answer{B}

\question Which is a KRI for crisis management readiness?
A) Loan balance
B) Time since last crisis simulation test
C) Deposit balance
D) Employee count
\answer{B}

\question What is "risk reporting anomaly detection"?
A) No detection
B) Identifying unusual patterns in risk metrics
C) Ignoring anomalies
D) Only manual review
\answer{B}

\question Which is a KRI for regulatory exam findings?
A) Loan growth
B) Number of repeat or overdue regulatory findings
C) Deposit balance
D) Employee count
\answer{B}

\question What is "risk monitoring escalation trigger"?
A) No trigger
B) Specific metric value that mandates escalation
C) Random trigger
D) Subjective trigger
\answer{B}

\question Which is a KRI for insurance coverage risk?
A) Loan balance
B) Percentage of uninsured or underinsured exposures
C) Deposit balance
D) Employee count
\answer{B}

\question What is "risk reporting quality review"?
A) No review
B) Periodic check of risk report accuracy and timeliness
C) No quality check
D) Only at year-end
\answer{B}

% ============================================================
% FINAL QUESTIONS 771-800 (Mixed Domains)
% ============================================================

\subsection{Comprehensive Review (Questions 771-800)}

\question What is the difference between "risk appetite" and "risk capacity"?
A) Same
B) Appetite is what you want to take; capacity is what you can take
C) Capacity is what you want
D) Appetite is larger
\answer{B}

\question Which is a key component of a risk management policy?
A) No scope
B) Scope, roles, requirements, and exception process
C) Marketing content only
D) Employee benefits
\answer{B}

\question What is "operational risk event" categorization?
A) No categorization
B) Internal fraud, external fraud, systems failures, process management
C) Only credit events
D) Only market events
\answer{B}

\question Which is a key output of an issue management process?
A) No output
B) Remediated issues with validated closure evidence
C) Unresolved issues only
D) Ignored findings
\answer{B}

\question What is the purpose of a "risk control matrix"?
A) No purpose
B) Mapping risks to controls for gap identification
C) Marketing tool
D) HR tool
\answer{B}

\question Which is a regulatory expectation for recovery planning?
A) No plan
B) Recovery plan for large banks to restore financial strength
C) Only for small banks
D) Optional
\answer{B}

\question What is "risk appetite misalignment"?
A) No issue
B) When business decisions are inconsistent with approved risk appetite
C) Perfect alignment
D) Intentional deviation
\answer{B}

\question Which is a key element of a risk report?
A) No context
B) Context, metrics, breaches, trends, and actions
C) Only numbers
D) Only text
\answer{B}

\question What is the purpose of "risk training for board members"?
A) No purpose
B) To ensure directors understand risk concepts and their oversight role
C) To replace management
D) To eliminate risk
\answer{B}

\question Which is a key measure of risk management effectiveness?
A) Number of policies
B) Risk incidents within appetite, issue closure rates, audit results
C) Number of employees
D) Number of branches
\answer{B}

\question What is "risk culture survey"?
A) No survey
B) Tool to assess employee perceptions of risk behaviors
C) Marketing survey
D) Customer survey
\answer{B}

\question Which is a regulatory expectation for third-party risk management?
A) No oversight
B) Inventory, due diligence, ongoing monitoring, and contingency plans
C) Ignore vendors
D) No contracts
\answer{B}

\question What is "risk appetite communication"?
A) No communication
B) Cascading risk appetite throughout the organization
C) Only to board
D) Only to regulators
\answer{B}

\question Which is a key output of risk and control self-assessment (RCSA)?
A) Marketing plan
B) Residual risk ratings and action plans
C) Employee schedule
D) Budget
\answer{B}

\question What is the purpose of "risk appetite metrics cascading"?
A) No purpose
B) Translating high-level appetite into business-level limits
C) Ignoring limits
D) Only at board level
\answer{B}

\question Which is a key principle of the COSO framework?
A) No monitoring
B) Risk assessment, control activities, information and communication
C) Only control activities
D) Only risk assessment
\answer{B}

\question What is "risk tolerance"?
A) Same as risk appetite
B) Specific quantitative boundaries within risk appetite
C) Only qualitative
D) Not related to appetite
\answer{B}

\question Which is a key output of stress testing?
A) Exact predictions
B) Identification of capital and liquidity needs under stress
C) No insights
D) Marketing insights
\answer{B}

\question What is the purpose of "risk policy exceptions"?
A) To avoid all rules
B) To allow necessary deviations with proper governance
C) To ignore policies
D) To eliminate policies
\answer{B}

\question Which is a key element of a risk management framework?
A) No governance
B) Governance, risk assessment, response, monitoring, reporting
C) Only monitoring
D) Only assessment
\answer{B}

\question What is "risk data aggregation"?
A) No aggregation
B) Combining risk data across systems for enterprise view
C) Siloed data only
D) Manual spreadsheets only
\answer{B}

\question Which is a regulatory expectation for internal audit?
A) No independence
B) Independence, competence, and direct board reporting
C) Part of management
D) No reporting
\answer{B}

\question What is the purpose of "risk limit structure"?
A) No purpose
B) To translate risk appetite into actionable boundaries
C) To eliminate all activity
D) To increase risk
\answer{B}

\question Which is a key output of scenario analysis?
A) No output
B) Vulnerability identification and contingency plans
C) Exact losses
D) No action
\answer{B}

\question What is "risk management maturity"?
A) No measurement
B) Assessment of risk capabilities compared to best practices
C) Only compliance
D) Only policies
\answer{B}

\question Which is a key element of a risk register?
A) No owner
B) Risk description, owner, rating, action plan, status
C) Only description
D) Only rating
\answer{B}

\question What is the purpose of "risk reporting to stakeholders"?
A) No purpose
B) To provide transparency and inform decision-making
C) To hide information
D) To confuse
\answer{B}

\question Which is a key measure of control effectiveness?
A) Control exists
B) Control operates as designed to reduce risk
C) Control is documented
D) Control has an owner
\answer{B}

\question What is "risk appetite statement review frequency"?
A) Once only
B) At least annually and when circumstances change
C) Never
D) Only after loss
\answer{B}

\question Which is a key principle of the three lines of defense?
A) No separation
B) First line owns risk, second line oversees, third line assures
C) First line audits
D) Second line owns risk
\answer{B}

% ============================================================
% ANSWER KEY (Questions 601-800)
% ============================================================
\newpage
\section*{Answer Key - Questions 601-800}
\begin{multicols}{3}
\small
\begin{enumerate}[start=601]
\item B \item B \item B \item B \item B \item B \item B \item B \item B \item B
\item B \item B \item B \item B \item B \item B \item B \item B \item B \item B
\item B \item B \item B \item B \item B \item B \item B \item B \item B \item B
\item B \item B \item B \item B \item B \item B \item B \item B \item B \item B
\item B \item B \item B \item B \item B \item B \item B \item B \item B \item B
\item B \item B \item B \item B \item B \item B \item B \item B \item B \item B
\item B \item B \item B \item B \item B \item B \item B \item B \item B \item B
\item B \item B \item B \item B \item B \item B \item B \item B \item B \item B
\item B \item B \item B \item B \item B \item B \item B \item B \item B \item B
\item B \item B \item B \item B \item B \item B \item B \item B \item B \item B
\item B \item B \item B \item B \item B \item B \item B \item B \item B \item B
\item B \item B \item B \item B \item B \item B \item B \item B \item B \item B
\item B \item B \item B \item B \item B \item B \item B \item B \item B \item B
\item B \item B \item B \item B \item B \item B \item B \item B \item B \item B
\item B \item B \item B \item B \item B \item B \item B \item B \item B \item B
\item B \item B \item B \item B \item B \item B \item B \item B \item B \item B
\item B \item B \item B \item B \item B \item B \item B \item B \item B \item B
\item B \item B \item B \item B \item B \item B \item B \item B \item B \item B
\item B \item B \item B \item B \item B \item B \item B \item B \item B \item B
\item B \item B \item B \item B \item B \item B \item B \item B \item B \item B
\end{enumerate}
\end{multicols}

\end{document}